\documentclass[12pt,a4paper,oneside]{article}
\usepackage[brazil]{babel}
\usepackage[utf8]{inputenc}
\usepackage{olymp}

\setcounter{problem}{0}

\begin{document}

\begin{problem}{Fatorial}{fatorial}{2}
 
O \textbf{fatorial} de um número natural $n$, representado por $n!$, é o produto de todos os inteiros positivos menores ou iguais a $n$. Ou seja,

$$n!=1 \times 2 \times 3 \times \dots \times (n-1) \times n$$

Além disso, por definição, $0! = 1$.

Faça um programa para, dado um número natural $n$, calcular seu fatorial $n!$.
%\Notes
%
%Observações, se tiver.

\InputFile

A entrada contém um único número natural $N$. Restrições: $0 \leq N \leq 20$.


\OutputFile

Seu programa deve gerar apenas uma linha de saída, contendo o valor de $N!$.

\Notes
Preste atenção aos tipos das variáveis, pois $20! \approx 2^{61}$.

\Example

\begin{example}
\exmp{
3
}{
6
}%
\end{example}

\begin{example}
\exmp{
5
}{
120
}%
\end{example}

\begin{example}
\exmp{
18
}{
6402373705728000
}%
\end{example}


%Comentários sobre o exemplo, se necessário.

\end{problem}

\end{document}
